% Formal benchmarks for the exact2021 cognitive presentation; no psychological adequacy claim.
\begin{annexitem}{annex:p19.r32.quantifiers}{thm:p19-r32-quantifier-benchmarks}{Universal application, existential limitation, and the pointwise alternative}
\begin{proof}
If a structure satisfies \(\forall x(Zx\to Gx)\) and \(Za\),
universal instantiation at the value of \(a\) gives \(Za\to Ga\),
and application gives \(Ga\). To refute the existential replacement,
use a two-element domain \(\{u,v\}\), interpret \(a\) by \(u\),
let \(Z\) hold of both elements and \(G\) only of \(v\).
The existential has witness \(v\) and \(Za\) holds, but \(Ga\)
is false. The constants/element names here are presentation data;
the premises contain no equality assumption.

For the pointwise alternative, fix an arbitrary element \(u\).
If it satisfies \(A\), the second universal premise gives \(Xu\).
The first gives \(Pu\lor Qu\), while the third gives \(\neg Qu\).
If the left disjunct holds, \(Pu\) is obtained directly; if the right
holds, it contradicts \(\neg Qu\), and the usual classical case
argument yields \(Pu\). Therefore \(Au\to Pu\) holds. Since
\(u\) was arbitrary, so does the universal conclusion. This uses
the pointwise negation premise, not an unanalyzed collective predicate.

For the weaker third premise take domain \(\{u,v\}\), make \(A,X\)
true of both, \(P\) true only of \(u\), and \(Q\) true only of
\(v\). Each element satisfies the first alternative and the
membership premise. The formula \(\forall x(Ax\to Qx)\) is false
at \(u\), so its negation is true; nevertheless the proposed
\(\forall x(Ax\to Px)\) is false at \(v\). This is the required
countermodel and makes the distinction between not-all and all-not
explicit.
\end{proof}
\end{annexitem}

\begin{annexitem}{annex:p19.r32.sentential}{thm:p19-r32-sentential-benchmarks}{The four-step chain and the material-conditional boundary}
\begin{proof}
Let a Boolean valuation satisfy all five displayed premises.
The alternative \(N\lor P\) and \(\neg N\) force \(P\).
The implication \(P\to H\) and that \(P\) force \(H\).
The implication \(H\to Z\) and that \(H\) force \(Z\).
If \(S\) were true, \(Z\land S\) would be true, contradicting
the last premise \(\neg(Z\land S)\). Thus \(\neg S\) follows.
This is the sequence \(P,H,Z,\neg S\), with each step explicitly
licensed by the declared Boolean meanings. Reversing premise order
changes none of their joint truth conditions, so the same sequence
is a possible derivation. Nothing in this argument predicts the
sequence actually reported by participants, its timing, or its
uniqueness among available valid proof routes.

Every valuation satisfying \(p\) satisfies \(p\lor q\), proving
that ordinary Boolean consequence. It does not require a particular
cognitive program to offer that conclusion to a participant.

For the conditional test, let the original premise set be
\(\{\neg q\}\). It continues to derive \(\neg q\) when \(p\)
is added as a supposition. But the original premises have a valuation
with \(p=0,q=0\), in which material \(p\to q\) is true.
Therefore deriving \(\neg q\) under the supposition does not imply
that the material conditional was false under the original information.
In contrast, the source's conditional-evaluation routine explicitly
assigns significance to the reasoning performed with the supposition;
that is its different psychological proposal, not a failure of this
truth-table calculation.
\end{proof}
\end{annexitem}

\begin{annexitem}{annex:p19.r32.beads}{thm:p19-r32-bead-sampling}{Two selection laws with the same box contents}
\begin{proof}
In the uniform hundred-bead experiment each bead has probability
\(1/100\); eighty disjoint red-origin outcomes therefore have
total probability \(80/100=.8\). In the two-stage experiment the
red box is chosen with probability \(1/2\). Each of its eighty beads
then has conditional probability \(1/80\), so each has unconditional
probability \(1/160\). Their sum is \(80/160=1/2\).
Similarly the twenty blue beads each have probability \(1/40\)
and their sum is \(1/2\), verifying normalization. Both experiments
have precisely the announced box contents and always select a bead
from one of the boxes, but have different red-origin probabilities.
Thus an additional sampling law is logically necessary to infer
the .80 answer from those counts. These laws are stipulated examples;
neither is an estimate of a human participant's response.
\end{proof}
\end{annexitem}

\begin{annexitem}{annex:p19.r32.generalization}{thm:p19-r32-generalization}{Why examples need scope, and why the decimal claim has a general proof}
\begin{proof}
Here is a concrete two-model witness for the first assertion. Use
domain \(\{0,1,2\}\) and a ternary relation \(R\), with the same
names denoting those elements in both structures. Record just
\(R(0,0,0)\) and \(R(1,0,1)\). In the first structure let
\(R=\{(0,0,0),(1,0,1),(2,0,2)\}\); then
\(\forall x\exists y\exists zR(x,y,z)\) is true. In the second
let \(R=\{(0,0,0),(1,0,1)\}\); the same recorded examples hold,
but there is no witness for \(x=2\), and that universal sentence
is false. Hence those finite positive records do not determine the
universal assertion. No claim has been made about a representation
that also supplies an exhaustive domain list or retains its quantified
input: those would be additional information and could exclude this
second model. This limited formal observation is not a psychological
impossibility theorem about model-based reasoning.

For the arithmetic assertion, write a finite decimal numeral as
\(n=\sum_{i=0}^{k}d_i10^i\), where each digit lies between zero
and nine. Ending in zero means \(d_0=0\). Thus
\[
 n=10\sum_{i=1}^{k}d_i10^{i-1}
   =5\left(2\sum_{i=1}^{k}d_i10^{i-1}\right).
\]
The parenthesized factor is a natural number, so five divides \(n\).
For the numeral zero the sum is empty and the factor is zero, which
also proves divisibility. This establishes all finite decimal numerals
with the specified ending, not merely a selected finite sample.
\end{proof}
\end{annexitem}
